\documentclass[10pt,a4paper]{amsart}
\usepackage[margin=19mm]{geometry}
\usepackage{amsmath,amssymb,mathtools}
\usepackage[scr=rsfs,cal=euler]{mathalfa}
\usepackage{xfrac}
\usepackage[unicode,hidelinks,bookmarks=false]{hyperref}
\newcommand{\ie}{i.e.\ }
\newcommand{\cf}{cf.\ }
\newcommand{\resp}{respectively\ }
\newcommand{\Subsection}{\subsection}
\providecommand{\nobreakdash}{\nobreak\hbox{-}}
\setlength{\emergencystretch}{2em}
\multlinegap0pt
\usepackage{amssymb}
\usepackage{xcolor}
\usepackage{algorithm}\usepackage{algpseudocode}
\usepackage{tikz}
\newtheorem{proposition}{Proposition}
\usepackage{cleveref}
\crefname{proposition}{Proposition}{Propositions}
\newcommand{\FirstHelix}{{\tt HEL}}
\DeclareMathOperator{\NB}{NB}
\newcommand{\STM}{{\tt STM}}
\title{Making ends meet or just meeting at the ends? Assessing end-to-end distance in folded RNA sequences and other branched structures}
\author{Torin Greenwood, Christine Heitsch}
\date{Source: arXiv:2604.15599v1 (CC BY 4.0)}
\begin{document}
\maketitle

\noindent\emph{Source and licence.} Making ends meet or just meeting at the ends? Assessing end-to-end distance in folded RNA sequences and other branched structures. Version arXiv:2604.15599v1 (CC BY 4.0), distributed by arXiv under the Creative Commons Attribution 4.0 International licence, http://creativecommons.org/licenses/by/4.0/, as recorded in the arXiv licence metadata for this exact version and shown on the abstract page https://arxiv.org/abs/2604.15599v1. Full licence notice: \texttt{licenses/arxiv-2604.15599-CC-BY-4.0.txt}.\par\smallskip

\noindent\textit{Editorial scope.} Counted unit: the article's proposition prop:DyckMotzkinDepth (source0.tex lines 746--754). The article's complete original proof of it is delivered with it (source0.tex lines 756--808, one \texttt{proof} environment of 53 lines). No mathematics is written, repaired or re-derived: the statement and the proof are the article's own lines, in the article's own notation, and no retained line is substituted or re-flowed.  Retained uncounted as the source-local context the proof needs: lines 450--464.  One declared adaptation: exactly 1 ancillary strings inside the retained spans are omitted (image-inclusion commands and, where the source repeats a label, the repeated label definitions) and no other line differs from the source; the figure environments, with their placement, captions and labels, are kept, and the package records the omitted strings in fidelity receipts bound to the affected bodies.  No retained line contains a citation, so the wrapper carries no bibliography and no bibliography entry.  The retained lines contain the article's own diagram code, which the wrapper typesets with the standard tikz packages; no image file is delivered and the package declares no asset dependency.  Editorial formatting only, in addition: an AMS \texttt{amsart} preamble that restates the article's own declarations and macros used by the retained lines, namely the environments \texttt{proposition} on separate counters, together with the standard packages those lines need.  The retained environments print in this document as Proposition 1 in order of appearance, whereas the article prints the same items with the numbering of its own full document.  The complete native source of the article as distributed in the arXiv e-print for this exact version is bundled unchanged as \texttt{sources/198624/source0.tex}.
\begin{proposition} \label{prop:GF}
Let $F(z, \mathbf{u})$ be a multivariate GF where $[z^n\mathbf{u}^\mathbf{k}]F(z, \mathbf{u})$ counts the number of objects of size $n$ with parameters equal to $\mathbf{k} = (k_1, k_2, \ldots, k_d)$.  Suppose $F(z, \mathbf{u})$ is of the form
   
    \[
    F(z, \mathbf{u}) = \frac{-b(z, \mathbf{u}) \pm g(z, \mathbf{u}) \sqrt{c(z)}}{2a(z, \mathbf{u})}
    \]
    with everywhere analytic functions $a, b, c,$ and $g$. Define $p_{n, \mathbf{k}}$ to be the probability that an object of size $n$ has parameters equal to $\mathbf{k}$ as below:
    \[
    p_{n, \mathbf{k}} = \frac{[z^n\mathbf{u}^\mathbf{k}] F(z, \mathbf{u})}{[z^n]F(z, \mathbf{1})}.
    \]
    Additionally, suppose there exists a $\rho > 0$ such that $c(\rho) = 0$ and $c(z) \neq 0$ for any other value of $z$ with $|z| \leq \rho$.  Assume that $c'(\rho) \neq 0$, and that for $\mathbf{u} \neq 0$ in a neighborhood of zero: $g(\rho, \mathbf{u}) \neq 0, a(\rho, \mathbf{u}) \neq 0$, and there are no non-removable singularities of $F(z, \mathbf{u})$ with $|z| \leq \rho$ besides $z = \rho$.  Also, assume $F(z, \mathbf{1})$ has no non-removable singularities with $|z| \leq \rho$ besides $z = \rho$. Then,
    \[
    \lim_{n \to \infty} p_{n, \mathbf{k}} = [\mathbf{u}^\mathbf{k}] \left(\frac{g(\rho, \mathbf{u})}{g(\rho, \mathbf{1})} \cdot \frac{a(\rho, \mathbf{1})}{a(\rho, \mathbf{u})}\right).
    \]
\end{proposition}

\begin{proposition} \label{prop:DyckMotzkinDepth}
Consider the uniform distribution over all Dyck or Motzkin paths of length $n$ as $n \to \infty$.
\begin{enumerate}
    \item For Dyck paths, \FirstHelix\ tends to $1 + \NB(1, 3/4)$ with mean $4/3$ and variance $4/9$.
    \item For Motzkin paths, \FirstHelix\ tends to $1 + \NB(1, 8/9)$ with mean $9/8$ and variance $9/64$.
    \item For Motzkin paths, the number of helices in the first stem tends to $1 + \NB(1, 27/32)$ with mean $32/27$ and variance $160/729$.
    \item For Motzkin paths, \STM\ tends to $1 + \NB(1, 3/4)$ with mean $4/3$ and variance $4/9$.
\end{enumerate}
\end{proposition}

\begin{proof}
    Our goal is to find a recurrence relation for each GF that tracks the corresponding depth of the paths.

    Let $D(z,y)$ be the GF which tracks Dyck paths by length ($z$) and height ($y$), where length is again the number of pairs of parentheses.  In order to find the recurrence for $D$, will need to define an auxiliary GF $H(x, y)$, described momentarily.  We begin with the standard Dyck path recurrence relation, modified slightly to include this helper GF:
    \[
    D(z, y) = 1 + [zy H(z,y)] D(z,1) 
    \]
    Overall, this recurrence is motivated by decomposing a Dyck path according to its first return to the $x$-axis, as in \Cref{fig:Dyck}.  A Dyck path is either empty (contributing a $1$), or it can be broken into a pair of parentheses enclosing a Dyck path, $(zy)$, the Dyck path in between the parentheses $(H(z, y))$, followed by another Dyck path $(D(z, 1))$.

    \begin{figure}
        \begin{center}
        
        \caption{A decomposition of Dyck paths according to the first return to the line $y = 0$.  The Dyck path within the box adds to the length of the first helix only if its first and last position are paired.} \label{fig:Dyck}
        \end{center}
    \end{figure}
    Now, we explain why the terms are $H(z, y)$ and $D(z, 1)$: for $D(z, 1)$, this is the univariate generating function for all Dyck paths by length, which we include because the Dyck path at the end of the recurrence cannot contribute to the depth of the first mountain.  Then, we define $H(z,y)$ to count all Dyck paths by length $(z)$ and depth $(y)$, except Dyck paths with 2+ mountains (or with zero) are now counted by coefficients of $z^i y^0$.  This is necessary because the depth of a path with 2+ mountains will be $1$ after the path is enclosed by the set of parentheses in the recurrence, and should be counted in the overall generating function $D(z, y)$ by the coefficient of $z^{i + 1} y^{0 + 1}$.  In other words, $H(z, y)$ only assigns positive depth to Dyck paths whose first step is paired with its last step.  We obtain the recurrence,
    \[
    H(z,y) = (1-z) D(z,1) + zy H(z,y),
    \]
    where the $(1-z)D(z, 1)$ term counts Dyck paths whose first step is not paired with its last step, and the $zyH(z, y)$ term counts the rest.

    Solving this system reveals that the GF $D(z, y)$ is of the form required in Proposition~\ref{prop:GF} with $g(y) = y$, $a(z, y) = y(z^2 - 1)$, and $\rho = 1/4$.  Plugging into the proposition yields the limiting distribution encoded by $3y/(4-y)$, which is the PGF for the geometric random variable with parameter $p = 3/4$.
   
    Paralleling the changes for the Dyck path specification, let $M(z, y)$ be the bivariate GF enumerating Motzkin paths by length ($z$, now counting individual nucleotides) and by depth of the first left-to-right mountain $(y)$.  Also, let $H(z, y)$ be the equivalently defined helper function.  Then, the following specification holds:
    \begin{align*}
    M(z, y) &= 1 + zM + z^2yH(z, y)M(z, 1)\\
    H &= z^2yH(z, y) + (1-z^2)M(z, 1)
    \end{align*}
    Solving for $M(z, y)$ again reveals a GF fitting the requirements of Proposition~\ref{prop:GF}, this time with $g(y) = y$, $a(z, y) = z^2(yz^3 - yz^2 - z + 1)$, and $\rho = 1/3$.  The proposition gives the limiting GF $(8y/9)/(1-y/9)$,
    which again corresponds to a geometric random variable, this time with parameter $p = 8/9$.

    For the length and number of components of the stem in a Motzkin path, we develop two grammars that output all dot-bracket sequences exactly once but allow us to track additions to the stem.  First, for \STM, we have the start symbol $M_f$ corresponding to when the grammar has yet to start generating the first stem.  Then, $M_*$ is the symbol corresponding to the first stem being in the middle of production, and $M$ is the symbol for when the grammar has moved past producing the stem.  Additionally, we have a symbol $L$ that corresponds to adding any number of consecutive dots.  We obtain:
    \begin{align*}
        M_f &\to \emptyset\ \big|\ .M_f\ \big|\ (M_*)M\\
        M_* &\to L\ \big|\ L(M_*)L\ \big|\ L(M)L(M)M\\
        M &\to \emptyset \ \big|\ .M \ \big|\ (M)M\\
        L &\to \emptyset \ \big|\ .L
    \end{align*}
    Here, $\emptyset$ represents outputting nothing.  The symbols $M_f$ and $M$ both have production rules modeled after the standard recursion for Motzkin paths.  In contrast, $M_*$ has rules broken into whether the output has no mountains $(L)$, one mountain $(L(M_*)L)$, or more than one mountain, which breaks the stem.  This converts to the system of equations
    \begin{align*}
        M_f &= 1 + zM_f + z^2uM_*M,\\
        M_* &= L + z^2uL^2M_* + z^4L^2M^3\\
        M&=1 + zM + z^2M^2\\
        L &= 1+zL
    \end{align*}
    This is solved via Gr\"{o}bner bases in our {\tt SageMath} worksheet, leading to a generating function where Proposition~\ref{prop:GF} applies.  The limiting GF is $3u/(4-u)$, again corresponding to the PGF for a geometric random variable with $p = 3/4$.

    The grammar needs a slight modification to track the number of helices in the stem.  In the rules for $M_*$, we now need to know more information about the production $M_* \to L(M_*)L$: if both $L$s are empty, a pre-existing helix is growing longer, but if either $L$ is non-empty, a new helix is formed in the stem instead.  This change corresponds to replacing the equation for $M_*$ with
    \[
    M_* = L + z^2M_* + z^2u(2(L-1) + (L-1)^2)M_* + z^4L^2M^3.
    \]
    Every other part of the analysis is the same, and the limiting PGF is now $27u/(32 - 5u)$, corresponding to the geometric random variable with parameter $p = 27/32$.
\end{proof}

\end{document}
